\documentclass[preprint,12pt]{elsarticle}

\usepackage{amsmath}
\usepackage{amsthm}
\usepackage{amssymb}
\usepackage{array}
\usepackage{bm}
\usepackage{booktabs}
\usepackage{mathtools}

\setlength{\emergencystretch}{2em}

\journal{Computational Materials Science}

\begin{document}

\begin{frontmatter}

\title{Model-Class Expressiveness Beyond Spectral Fitting: A Gauge-Equivariant Operator Framework}

\author[dlsu]{Eugene Joseph M. Ragasa\corref{cor1}}
\ead{eugene.ragasa@dlsu.edu.ph}
\cortext[cor1]{Corresponding author.}

\address[dlsu]{Department of Physics, De La Salle University,
2401 Taft Avenue, Manila 0922, Philippines}

\begin{abstract}
First-principles electronic Hamiltonians may be reduced to compact lattice
models by preserving either selected spectral properties or an aligned
localized operator. These criteria need not admit the same reduced
representation. This study formulates their compatibility within a common
hierarchy of orthogonal $sp^3s^{\ast}$ Slater--Koster models for bulk silicon.
A validated ten-orbital Wannier Hamiltonian will be constructed from one
converged density-functional-theory calculation. For each tight-binding model
class, the spectral and operator criteria define admissible sets satisfying
declared error tolerances. Compatibility requires a nonempty intersection,
corresponding to a single Hamiltonian satisfying both requirements. When the
intersection is empty, a normalized real-space Hamiltonian metric will
determine the minimum separation between the admissible sets and quantify the
incompatibility of the prescribed model class with the retained
first-principles information. Validation will use withheld band energies, the
indirect gap, conduction-valley position, longitudinal and transverse electron
effective masses, and operator residuals resolved by orbital block and neighbor
shell. The study will identify the smallest tested model class admitting a
common reduced representation satisfying both criteria and thereby establish
a validated bulk reference for subsequent impurity-operator reduction.
\end{abstract}

\begin{keyword}
silicon \sep Wannier Hamiltonian \sep tight-binding reduction
\sep operator compatibility \sep model selection
\end{keyword}

\end{frontmatter}

\section{Introduction}
\label{sec:introduction}

Density-functional theory (DFT) supplies a first-principles Kohn--Sham
description of crystalline electronic structure, but its plane-wave
representation is inconvenient for interpolation and reduced lattice
calculations. Maximally localized Wannier functions provide a localized
representation of a selected Kohn--Sham subspace and support accurate
interpolation~\cite{marzari2012,pizzi2020,souza2001,yates2007}. Parameterized
tight-binding models instead restrict the Hamiltonian to a prescribed orbital
basis, hopping range, and symmetry structure~\cite{slater1954,vogl1983}.

A spectral fit preserves selected eigenvalues, whereas an operator fit
preserves matrix information in an aligned localized representation. Spectral
agreement alone does not establish agreement of onsite terms, orbital
couplings, or hopping amplitudes. These differences become consequential when
the bulk Hamiltonian is later modified by dopants, defects, strain, or other
local perturbations, since all such operators must ultimately be represented
in the same reduced basis to form a consistent effective Hamiltonian. Before
an operator residual can be evaluated, the Wannier and tight-binding orbital
coordinates must therefore be identified through a stable, symmetry-compatible
alignment~\cite{bjork1973,higham1986}.

The present study asks whether the spectral and operator criteria admit a
common reduced representation within a prescribed tight-binding model class.
Rather than forcing two independently fitted parameter vectors to agree, it
defines admissible model sets for each criterion and asks whether they
intersect. Repeating this test over a nested hierarchy of tight-binding model
classes identifies the smallest class admitting a common representation, or
quantifies the obstruction when no compatible representation exists.

\section{Methodology}
\label{sec:methodology}

\subsection{First-principles reference calculation}
\label{sec:dft-reference}

\begin{sloppypar}
The first-principles reference will be generated using Quantum
ESPRESSO \cite{giannozzi2009}. Exchange and correlation will be treated using
the Perdew--Burke--Ernzerhof (PBE) functional \cite{perdew1996}.
An optimized norm-conserving pseudopotential from PseudoDojo
\cite{hamann2013,vansetten2018} will be used in scalar-relativistic form. The
parent problem will use the PBE-relaxed lattice constant and will be
non-spin-polarized without spin--orbit coupling. Plane-wave cutoffs, wavevector
meshes, and the number of bands will be converged against the indirect gap,
conduction-valley position, and longitudinal and transverse electron effective
masses.
\end{sloppypar}

\subsection{Wannier representation}
\label{sec:wannier-representation}

A ten-orbital target subspace containing the valence bands and the low
conduction states required for the silicon valleys will then be constructed
with Wannier90~\cite{pizzi2020,souza2001}. Its real-space Hamiltonian is

\begin{equation}
\left[
\bm{H}_{\mathrm{W}}(\bm{R})
\right]_{\alpha\beta}
=
\left\langle
w_{\alpha\bm{0}}
\middle|
\hat{H}_{\mathrm{KS}}
\middle|
w_{\beta\bm{R}}
\right\rangle,
\label{eq:wannier-hamiltonian}
\end{equation}

where $\bm{R}$ is a direct-lattice translation; $\alpha$ and $\beta$ label
Wannier orbitals in the reference and translated primitive cells,
respectively; $\lvert w_{\alpha\bm{0}}\rangle$ and
$\lvert w_{\beta\bm{R}}\rangle$ are localized Wannier states; and
$\hat{H}_{\mathrm{KS}}$ is the converged Kohn--Sham operator.

Initial projections and outer and frozen disentanglement windows will be
recorded. The Wannier Hamiltonian will be accepted as the common parent only
after its dense-grid interpolation error, centers, spreads, sensitivity to
window and projection choices, and real-space hopping decay have been
validated.

\subsection{Spectral and operator losses}
\label{sec:losses}

For each candidate model class, the spectral loss is a weighted error over
training-set Wannier eigenvalues and selected band-edge quantities. The
operator loss is

\begin{equation}
\mathcal{L}_{H}(\bm{C},\bm{\theta})
=
\sum_{\bm{R}}
\omega_{\bm{R}}
\left\|
\bm{H}_{\mathrm{W}}(\bm{R})
-
\bm{C}
\bm{H}_{\mathrm{TB}}(\bm{R};\bm{\theta})
\bm{C}^{\dagger}
\right\|_{\mathrm{F}}^{2},
\label{eq:operator-loss}
\end{equation}

where $\bm{C}\in\mathcal{U}$ is a symmetry-compatible unitary alignment from
the tight-binding orbital coordinates to the Wannier coordinates;
$\bm{\theta}$ is the tight-binding parameter vector;
$\bm{H}_{\mathrm{TB}}(\bm{R};\bm{\theta})$ is the tight-binding hopping
matrix associated with $\bm{R}$; $\omega_{\bm{R}}\geq 0$ weights the
contribution of each retained neighbor shell; and $\lVert\cdot\rVert_{\mathrm F}$
is the Frobenius norm.

\subsection{Tight-binding model hierarchy}
\label{sec:model-hierarchy}

The model hierarchy will begin with an orthogonal ten-orbital
$sp^3s^{\ast}$ Hamiltonian with nearest-neighbor hopping. If that class is
incompatible with the declared tolerances, second-neighbor and selected
symmetry-allowed corrections will be introduced in a fixed sequence.

For a candidate model class $\mathfrak{M}_{j}$, the two losses define the
admissible sets

\begin{align}
\mathfrak{A}_{E,j}(\tau_E)
&=
\left\{
\bm{H}(\bm{\theta})\in\mathfrak{M}_{j}
:
\mathcal{L}_{E}(\bm{\theta})\leq\tau_E
\right\},
\label{eq:spectral-admissible-set}
\\
\mathfrak{A}_{H,j}(\tau_H)
&=
\left\{
\bm{H}(\bm{\theta})\in\mathfrak{M}_{j}
:
\min_{\bm{C}\in\mathcal{U}}
\mathcal{L}_{H}(\bm{C},\bm{\theta})
\leq\tau_H
\right\},
\label{eq:operator-admissible-set}
\end{align}

where $\mathcal{L}_{E}$ is the spectral loss; $\mathcal{L}_{H}$ is the aligned
operator loss; $\tau_E$ is the declared spectral tolerance; $\tau_H$ is the
declared operator tolerance; and $j$ indexes the nested tight-binding model
classes.

The spectral tolerance includes targets of $5\%$ for the indirect gap, $3\%$
of the $\Gamma$--$X$ distance for the conduction-valley position, and $8\%$
for each electron effective mass. Operator and alignment tolerances will be
normalized to the converged Wannier Hamiltonian and fixed before the
compatibility search.

\subsection{Compatibility metric}
\label{sec:compatibility-metric}

A dimensionless real-space Hamiltonian metric $d_{\mathrm{TB}}$ will compare
models in their common canonical tight-binding coordinates. The minimum
distance between the two admissible sets is

\begin{equation}
\delta_{j}^{\ast}
=
\inf_{
\substack{
\bm{H}_{E}\in\mathfrak{A}_{E,j}\\
\bm{H}_{H}\in\mathfrak{A}_{H,j}
}
}
d_{\mathrm{TB}}(\bm{H}_{E},\bm{H}_{H}),
\label{eq:set-distance}
\end{equation}

where $\bm{H}_{E}$ and $\bm{H}_{H}$ are models satisfying the spectral and
operator criteria, respectively. A value of $\delta_{j}^{\ast}$ below the
declared map tolerance establishes compatibility within $\mathfrak{M}_{j}$;
otherwise, it quantifies the irreducible separation under the imposed spectral
and operator requirements.

All accepted models will be tested on withheld wavevectors. The operator
residual will be decomposed by onsite and hopping contribution, orbital block,
crystal-symmetry channel, and neighbor shell. The analysis is restricted to
bulk silicon; impurity operators are outside the present study.

All spectral and operator errors are measured relative to the converged
PBE/Wannier parent Hamiltonian. The present study tests reduction fidelity and
does not treat agreement with the experimental silicon band gap as an
independent fitting target.

\section{Expected results and decision criteria}
\label{sec:expected-results}

The required outputs are one converged bulk-silicon DFT dataset, one validated
ten-orbital Wannier Hamiltonian, admissible spectral and operator model sets for
each tested tight-binding class, their minimum metric separation, and one
common validation record. Numerical entries will be reported only after the
convergence, fitting, and sensitivity calculations are complete.

\begin{table}[t]
\centering
\caption{Decision criteria for the compatibility analysis.}
\label{tab:decision-criteria}
\begin{tabular}{@{}
>{\raggedright\arraybackslash}p{0.22\textwidth}
>{\raggedright\arraybackslash}p{0.29\textwidth}
>{\raggedright\arraybackslash}p{0.38\textwidth}@{}}
\toprule
Quantity & Mathematical criterion & Decision role \\
\midrule
Spectral admissibility
& Spectral loss within threshold
& Retains selected bands and band-edge quantities \\
Operator admissibility
& Aligned operator loss within threshold
& Retains real-space matrix information \\
Map compatibility
& Set distance within map threshold
& Establishes a common reduced Hamiltonian \\
Model selection
& Lowest-order compatible class
& Selects the minimal common Hamiltonian class \\
\bottomrule
\end{tabular}
\end{table}

The decisive result is the smallest tested model class for which the spectral
and operator admissible sets intersect within tolerance. If no tested class is
compatible, the nonzero set separation and its decomposition will identify the
orbital or hopping content responsible for the obstruction. Either outcome
provides a controlled basis for determining whether the chosen reduced
representation can consistently support subsequent dopant, defect, or impurity
operators.

This work addresses compatibility for a single pair of reduction criteria. The
same framework may later be evaluated for additional operator families within
separately authorized scope, where the objective is to determine whether a
common reduced representation exists for the bulk Hamiltonian together with
perturbation operators describing substitutional dopants or defects.

\section*{CRediT authorship contribution statement}

\textit{Working draft; confirm this statement with both authors before submission.}

Bryan D. Llenarizas: Formal analysis, Investigation, Validation, Data curation,
Visualization, Writing---original draft, Writing---review and editing.
Eugene Joseph M. Ragasa: Conceptualization, Methodology, Software, Validation,
Supervision, Project administration, Writing---review and editing.

\section*{Declaration of competing interest}

\textit{Working draft; confirm this declaration before submission.}

The authors declare that they have no known competing financial interests or
personal relationships that could have appeared to influence the work reported
in this paper.

\section*{Data and code availability}

The versioned source code, computational inputs, processed datasets, and
reproduction workflows supporting the final article will be deposited in
persistent repositories. Their permanent identifiers will be inserted here
before submission.

\section*{Declaration of generative AI and AI-assisted technologies}

A final declaration identifying the actual tools, uses, author review, and
responsibility statement will be completed before submission.

\section*{Acknowledgments}

% Add funding, computing-resource, institutional, and technical acknowledgments.
To be completed before submission.

\bibliographystyle{elsarticle-num}
\bibliography{references}

\end{document}
